% Einheit 3 von 4 – Neigungswinkel von Ebenen (bank/skalarprodukt-und-winkel/e3.jsonl, zusammenbau v0.1)
%% TODO Zeitmarke Einheit 3: Marken-Zeile nicht lesbar
%% TODO Zweigzeile Teil 1: was hier gelernt wird (Satz in Schülersprache aus dem Einheitstitel) – Einheit 3
\einheitenkopf[e3]{Einheit 3 von 4 · Neigungswinkel von Ebenen}
\zweigzeile{Abitur GK}
\setcounter{aufgabe}{13}

%% TODO Titel: Ich-kann-Satz fehlt in der Bank (Platzhalter: Kettenname) – Nr. 14
%% TODO Anweisung: ein Satz über der ersten Teilaufgabe fehlt in der Bank – Nr. 14
\begin{aufgabe}{Neigungswinkel von Ebenen}
%% TODO \rechenplatz{n}: Schrittzahl des Grundfalls fehlt in der Bank (nur bei fester Schreibform, 2.3 b)
\begin{teile}
\teil Der Winkel zwischen zwei Ebenen wird über ihre Normalenvektoren berechnet. Welche Funktion steht in der Formel? \\ \kreuz{Kosinus} \\ \kreuz{Sinus}
\teil Berechne den Winkel zwischen der Ebene $E\colon x + 2y + 2z = 8$ und der xy-Ebene. $\varphi \approx$ \leerfeld[$^\circ$]
\teil Zwei Seitenflächen eines Körpers liegen in den Ebenen $E\colon 2x + y + 2z = 10$ und $F\colon x - 2y + 2z = 3$. Berechne den Schnittwinkel der beiden Ebenen. $\varphi \approx$ \leerfeld[$^\circ$]
\teil Eine Kletterwand liegt in der Ebene $L\colon 6x + 3y + 2z = 24$, der Untergrund ist die xy-Ebene; 1 LE entspricht 1 m. Berechne die Größe des Winkels zwischen Kletterwand und Untergrund. \hfill (Abitur 2018 GK) \\ $\varphi \approx$ \leerfeld[$^\circ$]
\teil Eine Rollstuhlrampe liegt in der Ebene $x + 20z = 20$, der Boden ist die xy-Ebene. Ihre Neigung darf höchstens $6\,\%$ betragen. Prüfe, ob die Rampe diese Bedingung erfüllt.
\teil Eine Kellerklappe ist an der x-Achse drehbar befestigt und steht geöffnet in der Ebene $2y - 3z = 0$; der Boden ist die xy-Ebene. Beim Schließen wird die Klappe gedreht, bis der keilförmige Raum unter ihr kein Volumen mehr hat. Berechne den Drehwinkel und erläutere den Ansatz. \hfill (Abitur 2026 GK)
\end{teile}
\end{aufgabe}

%% TODO Titel: Ich-kann-Satz fehlt in der Bank (Platzhalter: Kettenname) – Nr. 15
%% TODO Anweisung: ein Satz über der ersten Teilaufgabe fehlt in der Bank – Nr. 15
\begin{aufgabe}{Neigungswinkel von Ebenen – Fehler finden, Begründen, Anwendung}
\begin{teile}
\teil Ben berechnet den Winkel zwischen der Ebene $2x + 6y + 3z = 7$ und der xy-Ebene: $\mathrm{sin}\,\varphi = \frac{3}{7}$, $\varphi \approx 25{,}4^\circ$. Finde den Fehler.
\teil Begründe, warum der Winkel zwischen den Normalenvektoren zweier Ebenen ebenso groß ist wie der Winkel zwischen den Ebenen.
\teil Ein Solarmodul liegt in der Ebene $2x + y + 3z = 15$, die xy-Ebene ist die Horizontale. Am meisten Strom liefert es bei einer Neigung zwischen $30^\circ$ und $40^\circ$. Liegt die Neigung in diesem Bereich?
\end{teile}
\end{aufgabe}
